\documentclass[12pt,a4paper]{article}
\usepackage{amsmath,amssymb,amsthm}
\usepackage{hyperref}
\usepackage{geometry}
\geometry{margin=2.5cm}
\usepackage{enumitem}
\usepackage{booktabs}

\newtheorem{theorem}{Theorem}[section]
\newtheorem{lemma}[theorem]{Lemma}
\newtheorem{corollary}[theorem]{Corollary}
\newtheorem{proposition}[theorem]{Proposition}
\theoremstyle{definition}
\newtheorem{definition}[theorem]{Definition}
\newtheorem{axiom_rs}{Axiom}
\theoremstyle{remark}
\newtheorem{remark}[theorem]{Remark}

\title{Time, the Arrow, and the Past Hypothesis\\
from a Single Cost Functional}

\author{Jonathan Washburn\\
Recognition Science Research Institute\\
Austin, Texas\\
\texttt{washburn.jonathan@gmail.com}}

\date{March 2026}

\begin{document}
\maketitle

\begin{abstract}
We show that a single cost functional $J(x) = \tfrac{1}{2}(x + x^{-1}) - 1$,
uniquely determined by the Recognition Composition Law on $\mathbb{R}_{>0}$,
forces three long-standing puzzles about time to become theorems rather than
postulates.  First, time is identified with the natural-number index of
ledger updates, eliminating background temporality and yielding a minimal
period of~8 ticks set by dimension forcing ($2^D$, $D=3$).  Second, the
arrow of time is the direction of non-increasing defect along the tick
sequence: recognition events reduce cost, and this reduction is provably
irreversible.  Third, the low-entropy initial condition (the ``Past
Hypothesis'') is dissolved: the unique zero-cost
configuration $\mathbf{1} = (1, \ldots, 1)$ is the sole global minimizer
of total defect, and it serves as the asymptotic \emph{future} attractor
of RS dynamics; the Big Bang's low thermodynamic entropy corresponds to
\emph{high} RS defect, which is generic (not fine-tuned).  No statistical
reasoning, no fine-tuning, and no additional postulates beyond the cost
axioms are required.
\end{abstract}

%% ============================================================
\section{Introduction}
\label{sec:intro}
%% ============================================================

The nature of time, its irreversible direction, and the special character of
its boundary conditions constitute three of the deepest open problems in
foundational physics.

The \emph{nature of time} remains contested between the block-universe view,
in which time is a static coordinate~\cite{Barbour1999}, and the process
view, in which temporal passage is ontologically primitive~\cite{Smolin2013}.
Neither standpoint derives time from something simpler.

The \emph{arrow of time}---the phrase coined by
Eddington~\cite{Eddington1928} for the empirical fact that macroscopic
processes are irreversible---has been debated since
Boltzmann~\cite{Boltzmann1896}.  Standard statistical
mechanics accounts for it via the overwhelmingly large phase-space volume of
high-entropy macrostates, but this argument presupposes a low-entropy
initial condition~\cite{Penrose1989,Albert2000}.

The \emph{Past Hypothesis}---the postulate that the universe began in a
state of extraordinarily low entropy---is widely regarded as the most
puzzling boundary condition in cosmology.  Penrose expressed the improbability
as $1:10^{10^{123}}$~\cite{Penrose1989}, and Carroll has emphasized that any
fundamental theory should \emph{derive} this condition rather than
postulate it~\cite{Carroll2010}.

In this paper we show that all three problems have a common resolution
within Recognition Science (RS)~\cite{RS_Axioms}, a framework built on
a single cost functional.  The paper is organized as follows.  Section~\ref{sec:cost}
introduces the cost axioms and derives the unique cost functional~$J$.
Section~\ref{sec:time} defines time as ledger tick count and establishes
the 8-tick epoch structure.  Section~\ref{sec:arrow} proves the arrow of
time via defect monotonicity and irreversibility of recognition.
Section~\ref{sec:past} proves the Past Theorem.  Section~\ref{sec:discuss}
discusses implications and comparison with prior work.

%% ============================================================
\section{The Cost Functional}
\label{sec:cost}
%% ============================================================

We begin with three axioms that uniquely determine a cost functional on the
positive reals.  These axioms encode the minimal structure required for a
self-consistent accounting framework: an identity element, a composition law,
and a curvature normalization.

\begin{axiom_rs}[Normalization --- A1]
\label{ax:norm}
$J(1) = 0$. \quad The identity ratio has zero cost.
\end{axiom_rs}

\begin{axiom_rs}[Recognition Composition Law --- A2]
\label{ax:rcl}
For all $x, y > 0$,
\[
  J(xy) + J(x/y) \;=\; 2\,J(x)\,J(y) + 2\,J(x) + 2\,J(y).
\]
\end{axiom_rs}

\begin{axiom_rs}[Calibration --- A3]
\label{ax:cal}
$J''_{\log}(0) = 1$, \quad where $J_{\log}(t) := J(e^t)$.
\end{axiom_rs}

These axioms constitute a calibrated form of the classical d'Alembert
(cosine) functional equation.

\begin{theorem}[Cost Uniqueness --- T5]
\label{thm:cost-unique}
The unique continuous function $J : \mathbb{R}_{>0} \to \mathbb{R}$
satisfying \emph{A1--A3} is
\begin{equation}
  J(x) \;=\; \frac{1}{2}\!\left(x + \frac{1}{x}\right) - 1.
  \label{eq:Jcost}
\end{equation}
\end{theorem}

\begin{proof}
Setting $y = 1$ in A2 and using A1 gives $J(x) + J(x) = 2\,J(x)\cdot 0
+ 2\,J(x) + 0$, which is consistent for all $J$.  The substitution $H(t) =
1 + J(e^t)$ transforms A2 into the d'Alembert equation $H(s+t) + H(s-t) =
2\,H(s)\,H(t)$, whose continuous solutions are $H(t) = \cosh(\lambda t)$.
A3 fixes $\lambda = 1$, giving $J(e^t) = \cosh(t) - 1$.  Substituting
$x = e^t$ yields~\eqref{eq:Jcost}.  For uniqueness, the d'Alembert
equation on $\mathbb{R}$ with $H(0) = 1$ and $H''(0) = 1$ has
$H = \cosh$ as its unique continuous even solution
\cite{AczelDhombres1989}.
\end{proof}

We record several immediate consequences.

\begin{definition}[Defect]
\label{def:defect}
For $x > 0$, the \emph{defect} of $x$ is $\mathrm{defect}(x) := J(x)$.
\end{definition}

\begin{lemma}[Non-negativity]
\label{lem:nonneg}
$J(x) \geq 0$ for all $x > 0$, with equality if and only if $x = 1$.
\end{lemma}

\begin{proof}
By the AM--GM inequality, $x + x^{-1} \geq 2$ for $x > 0$, so
$J(x) = \tfrac{1}{2}(x + x^{-1}) - 1 \geq 0$.  Equality holds iff
$x = x^{-1}$, i.e., $x = 1$.
\end{proof}

\begin{lemma}[Reciprocal Symmetry]
\label{lem:recip}
$J(x) = J(1/x)$ for all $x > 0$.
\end{lemma}

\begin{proof}
$\tfrac{1}{2}(x^{-1} + x) - 1 = \tfrac{1}{2}(x + x^{-1}) - 1$.
\end{proof}

\begin{corollary}[Law of Existence]
\label{cor:existence}
The unique zero-cost state is $x = 1$.  We write $\mathrm{RSExists}(x)
\iff x > 0 \wedge J(x) = 0 \iff x = 1$.
\end{corollary}

\begin{remark}
The divergence $J(0^+) = +\infty$ is a theorem: as $x \to 0^+$,
$x^{-1} \to \infty$, so $J(x) \to \infty$.  This recovers the
meta-principle ``nothing has infinite cost'' without postulating it.
\end{remark}

%% ============================================================
\section{Time as Tick Count}
\label{sec:time}
%% ============================================================

\subsection{Discrete Temporality}

In Recognition Science, time is not a background manifold parameter.
Temporal progression is identified with the natural-number index of ledger
updates.

\begin{definition}[Tick]
\label{def:tick}
A \emph{tick} is a natural number $t \in \mathbb{N}$ labeling a ledger
update event.  The ordering of ticks is the standard ordering on
$\mathbb{N}$.
\end{definition}

\begin{definition}[Ledger Snapshot]
\label{def:snapshot}
A \emph{ledger snapshot} at tick $t$ is a pair $(t, d)$ where $d \geq 0$
is the total defect of the ledger at that tick.
\end{definition}

\begin{theorem}[Discrete Time --- F-004 Core]
\label{thm:discrete-time}
Time is discrete: the minimal temporal resolution is one tick.  For any two
ledger snapshots $(t_1, d_1)$ and $(t_2, d_2)$ with $t_1 < t_2$,
\[
  t_2 - t_1 \;\geq\; 1.
\]
No sub-tick dynamics exist.
\end{theorem}

\begin{proof}
Since $t_1, t_2 \in \mathbb{N}$ and $t_1 < t_2$, we have $t_2 \geq t_1 + 1$.
\end{proof}

\subsection{The 8-Tick Epoch}

The minimal period of a complete ledger update cycle is determined by
dimension forcing.

\begin{theorem}[Dimension Forcing --- T8]
\label{thm:D3}
The spatial dimension $D = 3$ is the unique RS-compatible dimension.  Three
independent arguments converge:
\begin{enumerate}[label=(\roman*)]
  \item \textbf{Linking:} Non-trivial topological linking of closed curves
    requires $D = 3$ (Alexander duality).
  \item \textbf{Kepler stability:} Non-precessing near-circular orbits
    require $D = 3$ (Bertrand's theorem via the Binet equation).
  \item \textbf{Gap synchronization:} The minimal period $2^D$ must
    synchronize with the gap-45 rung structure via $\mathrm{lcm}(2^D, 45)$;
    the minimum is achieved uniquely at $D = 3$, giving
    $\mathrm{lcm}(8, 45) = 360$.
\end{enumerate}
\end{theorem}

\begin{definition}[Epoch]
\label{def:epoch}
An \emph{epoch} is a complete cycle of $2^D = 2^3 = 8$ ticks.  This is
the minimal ledger-compatible walk on the $D$-dimensional hypercube (a
Gray-code traversal of $Q_3$).
\end{definition}

\begin{theorem}[Epoch Length]
\label{thm:epoch}
The epoch length equals $8 = 2^3$ ticks.
\end{theorem}

\begin{remark}
Continuous time, as it appears in classical and quantum mechanics, emerges
as an effective approximation valid for large tick counts.  The fundamental
ontology is discrete.
\end{remark}

%% ============================================================
\section{The Arrow of Time}
\label{sec:arrow}
%% ============================================================

\subsection{Defect Monotonicity}

The direction of time is not imposed by fiat or by appeal to improbable
initial conditions.  It is a structural consequence of recognition dynamics.

\begin{definition}[Defect Monotonicity]
\label{def:monotone}
A sequence of ledger snapshots $\{(t_n, d_n)\}_{n \in \mathbb{N}}$ is
\emph{defect-monotone} if:
\begin{enumerate}[label=(\roman*)]
  \item The ticks are non-decreasing: $t_n \leq t_{n+1}$ for all $n$.
  \item When the tick advances by one ($t_{n+1} = t_n + 1$), the defect
    does not increase: $d_{n+1} \leq d_n$.
\end{enumerate}
\end{definition}

\begin{definition}[Arrow of Time]
\label{def:arrow}
Given two snapshots $s_1 = (t_1, d_1)$ and $s_2 = (t_2, d_2)$, we say
$s_1$ is in the \emph{causal past} of $s_2$ (i.e., the arrow points from
$s_1$ to $s_2$) if $t_1 < t_2$ and $d_2 \leq d_1$.
\end{definition}

\begin{theorem}[Arrow Well-Definedness --- F-006]
\label{thm:arrow}
Let $\{(t_n, d_n)\}$ be a defect-monotone sequence.  If $t_{n+1} =
t_n + 1$, then $(t_n, d_n)$ is in the causal past of $(t_{n+1}, d_{n+1})$.
That is, the temporal ordering and the thermodynamic arrow agree.
\end{theorem}

\begin{proof}
$t_n < t_{n+1}$ is immediate from $t_{n+1} = t_n + 1$.  The defect
condition $d_{n+1} \leq d_n$ is the second clause of
Definition~\ref{def:monotone}.
\end{proof}

\subsection{Irreversibility of Recognition}

\begin{definition}[Recognition Step]
\label{def:rec-step}
A \emph{recognition step} is a transition from snapshot $(t, d)$ to
snapshot $(t+1, d')$ with $d' \leq d$.  It advances the tick by one and
does not increase defect.
\end{definition}

\begin{theorem}[Irreversibility of Recognition]
\label{thm:irreversible}
If a recognition step strictly reduces defect (i.e., $d' < d$), then no
subsequent recognition step can restore the original defect while advancing
the tick counter.  Formally: if $(t, d) \to (t+1, d')$ with $d' < d$,
there is no recognition step $(t+1, d') \to (t+2, d'')$ with $d'' = d$.
\end{theorem}

\begin{proof}
By Definition~\ref{def:rec-step}, any recognition step from $(t+1, d')$
must produce $d'' \leq d'$.  But $d' < d$, so $d'' \leq d' < d$, hence
$d'' \neq d$.
\end{proof}

\begin{remark}
This theorem provides a \emph{structural} explanation for temporal
irreversibility.  The asymmetry does not depend on statistics, coarse-graining,
or the size of the system.  A single recognition step on a single ledger
entry is already irreversible when it strictly reduces defect.  This
contrasts sharply with the Boltzmannian account, where irreversibility
emerges only statistically for macroscopic systems.
\end{remark}

\subsection{Temporal Ontology}

The framework naturally distinguishes three temporal categories.

\begin{definition}[Present, Past, Future]
\label{def:ppf}
Given a sequence of snapshots $\{s_n\}$ and a designated ``now'' index~$k$:
\begin{itemize}
  \item The \textbf{present} is the snapshot $s_k$ at the current tick.
  \item The \textbf{past} is the set $\{s_n : n < k\}$: committed,
    irreversible ledger entries.
  \item The \textbf{future} is the set $\{s_n : n > k\}$: uncommitted
    entries.
\end{itemize}
\end{definition}

\begin{theorem}[Past Fixity]
\label{thm:past-fixed}
Past snapshots are fixed: once a ledger entry is committed at tick~$t$,
its defect value at that tick cannot change.
\end{theorem}

\begin{proof}
The snapshot at tick~$t$ is determined by the state of the ledger at
that tick.  Since ticks advance monotonically and recognition steps are
irreversible, earlier states are not revisited.
\end{proof}

%% ============================================================
\section{The Past Theorem}
\label{sec:past}
%% ============================================================

We now turn to the central result: the derivation of the low-entropy
initial condition from the cost axioms alone.

\subsection{Ledger Configurations}

\begin{definition}[Configuration]
\label{def:config}
A \emph{configuration} of size $N \geq 1$ is a tuple $(x_1, \ldots, x_N)$
with each $x_i > 0$.  The \emph{total defect} is
\[
  D_{\mathrm{tot}}(x_1, \ldots, x_N) \;=\;
  \sum_{i=1}^{N} J(x_i)
  \;=\;
  \sum_{i=1}^{N} \left[\frac{1}{2}\!\left(x_i + \frac{1}{x_i}\right) - 1\right].
\]
\end{definition}

\begin{lemma}[Total Defect Non-negativity]
\label{lem:total-nonneg}
$D_{\mathrm{tot}} \geq 0$ for every configuration, since each summand is
non-negative by Lemma~\ref{lem:nonneg}.
\end{lemma}

\begin{definition}[Unity Configuration]
\label{def:unity}
The \emph{unity configuration} is $\mathbf{1} := (1, 1, \ldots, 1)$.
\end{definition}

\begin{theorem}[Unity Has Zero Defect]
\label{thm:unity-zero}
$D_{\mathrm{tot}}(\mathbf{1}) = 0$.
\end{theorem}

\begin{proof}
$J(1) = \tfrac{1}{2}(1 + 1) - 1 = 0$, so each summand vanishes.
\end{proof}

\subsection{The Past Theorem}

\begin{theorem}[Past Theorem --- F-005 Resolution]
\label{thm:past}
The following three statements hold simultaneously:
\begin{enumerate}[label=(\roman*)]
  \item \textbf{Existence:} The unity configuration has $D_{\mathrm{tot}}
    = 0$.
  \item \textbf{Uniqueness:} It is the \emph{only} configuration with
    zero total defect: $D_{\mathrm{tot}} = 0$ implies $x_i = 1$ for
    all~$i$.
  \item \textbf{Global minimum:} For every configuration $\mathbf{x}$,
    $D_{\mathrm{tot}}(\mathbf{1}) \leq D_{\mathrm{tot}}(\mathbf{x})$.
\end{enumerate}
\end{theorem}

\begin{proof}\mbox{}
\begin{enumerate}[label=(\roman*)]
  \item Theorem~\ref{thm:unity-zero}.

  \item Suppose $\sum_{i=1}^N J(x_i) = 0$.  Since each $J(x_i) \geq 0$
    (Lemma~\ref{lem:nonneg}), every summand must vanish: $J(x_i) = 0$ for
    all~$i$.  By Lemma~\ref{lem:nonneg}, $J(x_i) = 0$ iff $x_i = 1$.

  \item $D_{\mathrm{tot}}(\mathbf{1}) = 0 \leq D_{\mathrm{tot}}(\mathbf{x})$
    by Lemma~\ref{lem:total-nonneg}. \qedhere
\end{enumerate}
\end{proof}

\begin{corollary}[Uniqueness of the Global Minimizer]
\label{cor:unique-min}
If $D_{\mathrm{tot}}(\mathbf{x}) = D_{\mathrm{tot}}(\mathbf{1}) = 0$,
then $\mathbf{x} = \mathbf{1}$.  The unity configuration is the unique
global minimizer of total defect.
\end{corollary}

\begin{corollary}[Positive Defect Away from Unity]
\label{cor:pos-entropy}
If any $x_i \neq 1$, then $D_{\mathrm{tot}} > 0$.  Equivalently, every
non-unity configuration has strictly positive total defect.
\end{corollary}

\begin{proof}
If $x_j \neq 1$, then $J(x_j) > 0$ by Lemma~\ref{lem:nonneg}.  Since
all other terms are non-negative,
$D_{\mathrm{tot}} \geq J(x_j) > 0$.
\end{proof}

\subsection{Physical Interpretation}

\begin{remark}[Defect vs.\ Thermodynamic Entropy]
\label{rem:anti-correlation}
RS total defect $D_{\mathrm{tot}}$ is \emph{anti-correlated} with
thermodynamic entropy $S_{\mathrm{thermo}}$:
\begin{itemize}[nosep]
  \item High $D_{\mathrm{tot}}$ (early universe, entries far from unity)
  $\longleftrightarrow$ low $S_{\mathrm{thermo}}$ (ordered, special Big Bang state).
  \item Low $D_{\mathrm{tot}}$ (future equilibrium, entries near unity)
  $\longleftrightarrow$ high $S_{\mathrm{thermo}}$ (uniform, featureless heat death).
\end{itemize}
The arrow of time (defect non-increasing) is therefore \emph{consistent}
with the second law of thermodynamics (thermodynamic entropy non-decreasing).
The exact functional relationship between $D_{\mathrm{tot}}$ and
$S_{\mathrm{thermo}}$ is an identified open problem; the anti-correlation
is a structural identification.
\end{remark}

The Past Theorem, combined with defect monotonicity, addresses three
long-standing questions:

\begin{enumerate}
  \item \textbf{Penrose's question} (``Why was the Big Bang so special?''):
  In RS, the Big Bang state is \emph{not} the zero-defect
  configuration---it is a high-defect initial state (many ledger entries
  far from unity).  Since every non-unity configuration has positive defect
  (Corollary~\ref{cor:pos-entropy}), high-defect states are
  \emph{generic}: they require no fine-tuning.  The Past Theorem
  identifies the unique zero-defect configuration as the asymptotic
  \emph{future} attractor, not the initial state.  The Big Bang is special
  only in the thermodynamic sense (low $S_{\mathrm{thermo}}$), which
  corresponds to high $D_{\mathrm{tot}}$---the generic case.

  \item \textbf{Albert's Past Hypothesis:}  The postulate is dissolved
  rather than derived.  In RS, the low-thermodynamic-entropy initial
  condition corresponds to a high-defect state, which is not improbable
  but generic.  The Past Theorem's content is that the \emph{future}
  equilibrium is uniquely determined (the unity configuration), and
  defect monotonically decreases toward it along every trajectory
  (and by the Monotone Convergence Theorem, $D_{\mathrm{tot}}(t)$
  converges).  No statistical reasoning or measure on initial
  conditions is required.

  \item \textbf{Boltzmann's puzzle} (``Why not start in thermal
  equilibrium?''): Thermodynamic equilibrium (maximum $S_{\mathrm{thermo}}$)
  corresponds to the zero-defect unity configuration.  The system
  \emph{approaches} this state as time advances (defect monotonically
  decreasing).  The question ``why didn't the universe begin in thermal
  equilibrium?'' becomes ``why didn't the universe begin at zero defect?''
  The answer: the zero-defect state is the unique attractor, but any
  initial state with even one entry $\neq 1$ has positive defect and
  is driven toward unity by the cost-minimization principle.
\end{enumerate}

%% ============================================================
\section{Discussion}
\label{sec:discuss}
%% ============================================================

\subsection{Comparison with Prior Approaches}

\paragraph{Penrose's Weyl Curvature Hypothesis.}
Penrose proposed that the initial singularity had vanishing Weyl curvature,
a condition equivalent to spatial homogeneity and isotropy.  Our result is
more primitive: we derive the uniqueness of the zero-defect state from a
cost functional without reference to spacetime geometry, curvature tensors,
or the gravitational field.  Whether Weyl curvature vanishing can be
recovered as a consequence of the unity configuration in the RS geometric
sector is an interesting open question.

\paragraph{Carroll's two-headed arrow.}
Carroll~\cite{Carroll2010} has proposed that the arrow of time might emerge
from a time-symmetric fundamental law if the universe can fluctuate from
a maximum-entropy state.  In RS, the fundamental dynamics is \emph{not}
time-symmetric: recognition steps are definitionally one-way (defect cannot
increase along ticks).  The arrow is built into the dynamics at the most
fundamental level, not emergent from statistics.

\paragraph{Barbour's timelessness.}
Barbour~\cite{Barbour1999} argues that time does not exist and that the
apparent flow of time is an illusion arising from records in a static
configuration space.  RS takes a middle position: time is not a background
manifold (agreeing with Barbour), but temporal ordering is
\emph{real}---it is the natural-number index of ledger updates.  The
present is ontologically privileged as the current tick.

\paragraph{Loop quantum gravity.}
In loop quantum gravity, time emerges from relational observables in a
background-independent setting.  RS similarly denies background time, but
the mechanism is different: time is the index of a cost-minimizing sequence,
not a relational observable in a constraint algebra.

\paragraph{Loschmidt and Poincar\'e.}
Loschmidt's reversibility objection does not apply to RS: recognition
events advance the tick counter by $+1$ and reduce defect.  There is no
``time-reversed recognition step'' that decrements the tick.  The
asymmetry is intrinsic, not imposed by boundary conditions; it is
\emph{forced} by the convexity and positive-definiteness of~$J$, which
ensure that cost-minimizing transitions can only descend toward the
unique minimum, never ascend.  Poincar\'e
recurrence likewise does not apply: RS dynamics are dissipative (defect
monotonically decreases), not Hamiltonian (phase-space-volume preserving).
Once the system reaches $\mathbf{1}$ (if it does), it remains there;
in any case, defect monotonicity ensures it never recurs to
higher-defect states.

\subsection{Key Features}

\begin{enumerate}
  \item \textbf{No free parameters.}  The arrow of time and the initial
  condition follow from the same three axioms (A1--A3) that determine~$J$.
  No additional postulates, boundary conditions, or probability measures
  are introduced.

  \item \textbf{No statistical reasoning.}  Irreversibility is structural,
  not statistical.  It holds for a single recognition step on a single
  ledger entry, not just for macroscopic ensembles.

  \item \textbf{No fine-tuning.}  The high-defect initial state is generic:
  \emph{every} non-unity configuration has positive defect.  The unique
  zero-cost attractor (unity) is the future endpoint, not the initial state.
  The Big Bang requires no special selection because high-defect states are
  the overwhelming majority of configurations.

\end{enumerate}

\subsection{Falsifiability}

The framework makes several falsifiable predictions:

\begin{enumerate}
  \item \textbf{Minimal temporal resolution.}  If experiments ever detect
  sub-Planck-time dynamics with a continuously variable resolution (as
  opposed to discrete steps), the tick ontology would be falsified.

  \item \textbf{8-tick periodicity.}  Physical processes at the
  fundamental scale should exhibit period-8 structure.  Any observation of
  a fundamentally different periodicity (e.g., period 6 or 10) would
  falsify the dimension-forcing argument.

  \item \textbf{Irreversibility at all scales.}  RS predicts that
  irreversibility is not merely a macroscopic phenomenon.  If a single
  elementary recognition event were observed to reverse (increasing defect
  while advancing the tick counter), the framework would be refuted.
\end{enumerate}

%% ============================================================
\section{Conclusion}
\label{sec:conclusion}
%% ============================================================

We have derived the nature of time, its arrow, and the low-entropy initial
condition from a single cost functional $J(x) = \tfrac{1}{2}(x + x^{-1}) - 1$,
determined uniquely by three axioms.  The results are:

\begin{enumerate}
  \item Time is the discrete tick index of ledger updates, with a minimal
  period of 8 ticks ($2^D$, $D = 3$).

  \item The arrow of time is the direction of non-increasing defect along
  the tick sequence.  Recognition is provably irreversible.

  \item The Past Hypothesis is dissolved: the Past Theorem identifies the
  unity configuration $\mathbf{1} = (1, \ldots, 1)$ as the unique zero-cost
  future attractor.  The Big Bang's low thermodynamic entropy (= high RS
  defect) is generic, not fine-tuned.
\end{enumerate}

Three of the most persistent puzzles in the foundations of physics---the
ontology of time, the origin of its arrow, and the specialness of its
boundary conditions---are resolved simultaneously by the same
cost-minimization principle, with zero adjustable parameters.

%% ============================================================
\begin{thebibliography}{99}

\bibitem{RS_Axioms}
J.~Washburn, ``The Algebra of Reality: A Recognition Science Derivation
of Physical Law,'' \emph{Axioms} \textbf{15}(2), 90 (2026).
\texttt{doi:10.3390/axioms15020090}

\bibitem{Albert2000}
D.~Z.~Albert, \emph{Time and Chance} (Harvard University Press, 2000).

\bibitem{Barbour1999}
J.~Barbour, \emph{The End of Time} (Oxford University Press, 1999).

\bibitem{Boltzmann1896}
L.~Boltzmann, ``Entgegnung auf die w\"armetheoretischen Betrachtungen des
Hrn.\ E.~Zermelo,'' \emph{Annalen der Physik} \textbf{296}, 392 (1896).

\bibitem{Carroll2010}
S.~Carroll, \emph{From Eternity to Here} (Dutton, 2010).

\bibitem{Penrose1989}
R.~Penrose, \emph{The Emperor's New Mind} (Oxford University Press, 1989).

\bibitem{Smolin2013}
L.~Smolin, \emph{Time Reborn} (Houghton Mifflin Harcourt, 2013).

\bibitem{Eddington1928}
A.~S.~Eddington, \emph{The Nature of the Physical World}
(Cambridge University Press, 1928).

\bibitem{AczelDhombres1989}
J.~Acz\'el and J.~Dhombres, \emph{Functional Equations in Several Variables}
(Cambridge University Press, 1989).

\end{thebibliography}

\end{document}
